\section{Quantifier audit for logarithmic correlation theorems}
\label{sec:literature}

The two-point logarithmically averaged Chowla theorem controls fixed
affine forms on a multiplicative interval with the reciprocal weight
of the original arithmetic variable \citep{Tao2016LogChowla}.  The
averaged Chowla theorem averages the shifts
\citep{MatomakiRadziwillTao2015}.  Structural results at almost all
scales retain an almost-all qualifier \citep{TaoTeravainen2019}.
These are major arithmetic inputs, but their quantifiers are not
interchangeable with the literal TPC packet.

\begin{proposition}[Literal promotion checklist]
\label{prop:promotion-checklist}
A logarithmic theorem can be promoted to a certificate for
\eqref{eq:literal-reassembly} only after verifying:
\begin{enumerate}[label=\textup{(\roman*)},leftmargin=2.5em]
\item the same prescribed \(h_0\) and the actual slopes and
  intercepts, including their growth with \(X\);
\item the original arithmetic variable and therefore the true anchor
  \(\rho_{\theta,X}\);
\item uniformity in every prefix required by partial summation or
  dyadic recursion;
\item all interval origins, masks, smooth weights, periodic factors,
  polarizations, and native outer keys;
\item an exact return of short, resonant, content, and endpoint
  sectors; and
\item a quantitative aggregate rate meeting
  \eqref{eq:aggregate-target}, or a direct bound for the signed sum
  \eqref{eq:aggregate-signed}.
\end{enumerate}
\end{proposition}

\begin{proof}
Items (i)--(v) identify the sequence and frame in
\eqref{eq:fiber-anchored-primitive}.  Item (vi) is precisely the
hypothesis needed by \cref{thm:aggregate-identity}.  Omitting the
anchor contradicts \cref{prop:no-anchor-free}; omitting outer mass
contradicts \cref{ex:outer-mass}.
\end{proof}

\begin{remark}[Fixed forms versus growing forms]
Tao's theorem is stated for fixed \(a_1,a_2,b_1,b_2\) with nonzero
determinant.  The literal columns have slopes, intercepts, origins,
and masks depending on \(X\), while preserving one fixed nonzero
determinant \(h_0\) after the exact content reduction.  No uniformity
in that entire growing family is inferred here.
\end{remark}

\begin{remark}[Averaged shifts do not select \(h_0\)]
An estimate averaged over many shifts can show that most shifts are
good without identifying the one prescribed physical shift.  It
remains at its averaged level until a literal selection theorem is
proved.
\end{remark}

\begin{remark}[Logarithmic normalization]
A normalized bound \(L_{r,N}=o(\cH_{r,N})\) has different strength in
different regimes.  For \(r\ge N\), a prefix-uniform version is enough
by \cref{thm:short-prefix-uniform}.  For \(r=0\),
\(\cH_{0,N}\asymp\log N\), and qualitative
\(o(\log N)\) cancellation alone does not control the unshifted
signed defect; this is the sharp gap of TPC-87.
\end{remark}
